\section{Introduction}\label{sec:introduction}
The entropy photon number of an \(n\)-mode bosonic state is the mean photon number per mode of the product of \(n\) identical one-mode thermal states with the same total entropy. Guha, Erkmen, and Shapiro conjectured that this quantity obeys a concavity inequality under beam-splitter mixing of independent inputs \cite{GES2007}. Their entropy photon-number inequality is a quantum counterpart of the classical entropy power inequality and has consequences for minimum output entropy and the information capacities of bosonic channels. We prove its two-input form for finite-energy states.

Write
\begin{equation}\label{eq:g-definition}
 g(t)=(t+1)\log(t+1)-t\log t,\qquad t\ge0,
\end{equation}
with \(0\log0=0\). This continuous function increases from zero to infinity; its inverse is denoted by \(g^{-1}\). All entropies and logarithms in this paper use the natural base.

\begin{theorem}[Entropy photon-number inequality]\label{thm:epni}
Let \(n\ge1\) be finite, and let \(\rho_A,\rho_B\) be states on \(n\) bosonic modes, with annihilation operators \(a_1,\ldots,a_n\) and \(b_1,\ldots,b_n\), respectively, satisfying
\[
 \Tr\rho_A\sum_{j=1}^n a_j^\dagger a_j<\infty,
 \qquad
 \Tr\rho_B\sum_{j=1}^n b_j^\dagger b_j<\infty.
\]
The joint input is \(\rho_A\otimes\rho_B\); no independence assumption is made among the modes within either state. For \(0\le\eta\le1\), let \(\rho_C\) be the reduced state of the modes
\[
 c_j=\sqrt\eta\,a_j+\sqrt{1-\eta}\,b_j,\qquad 1\le j\le n,
\]
after passive beam-splitter mixing. Then
\begin{equation}\label{eq:epni}
 g^{-1}\!\left(\frac{S(\rho_C)}n\right)
 \ge \eta\,g^{-1}\!\left(\frac{S(\rho_A)}n\right)
 +(1-\eta)\,g^{-1}\!\left(\frac{S(\rho_B)}n\right),
\end{equation}
where \(S(\rho)=-\Tr\rho\log\rho\).
\end{theorem}

Thus Theorem~\ref{thm:epni} resolves the entropy photon-number conjecture positively for finite-energy inputs, including arbitrary internal multimode entanglement. Thermal inputs with a common mean photon number within each port give equality, even when the means at the two ports differ.

For the channel consequences, write \(\tau_r\) for the \(n\)-mode product thermal state with mean photon number \(r\) in each mode, including the vacuum at \(r=0\); thus \(S(\tau_r)=ng(r)\). Let \(\mathcal E_{\eta,N_B}\) be the one-mode thermal attenuator obtained by retaining \(c=\sqrt\eta\,a+\sqrt{1-\eta}\,b\) when the environment mode \(b\) is thermal with mean photon number \(N_B\) and is independent of the input. Its tensor power \(\mathcal E_{\eta,N_B}^{\otimes n}\) uses the same parameters in every mode and the joint input \(\rho\otimes\tau_{N_B}\).

\begin{corollary}[Constrained entropy minimum for identical thermal attenuators]\label{cor:thermal-attenuator}
Let \(n\ge1\) be finite, let \(s,N_B\ge0\) be finite, and let \(0\le\eta\le1\). Every state \(\rho\) on \(n\) modes with
\(\Tr\rho\sum_{j=1}^n a_j^\dagger a_j<\infty\) satisfies
\begin{equation}\label{eq:thermal-attenuator-bound}
 \frac1n S\!\left(\mathcal E_{\eta,N_B}^{\otimes n}(\rho)\right)
 \ge g\!\left(\eta g^{-1}\!\left(\frac{S(\rho)}n\right)
                   +(1-\eta)N_B\right).
\end{equation}
In particular, the exact constrained minimum is
\begin{equation}\label{eq:thermal-attenuator-minimum}
 \min_{\substack{\rho\text{ a state on }n\text{ modes}\\
                  \Tr\rho\sum_{j=1}^n a_j^\dagger a_j<\infty\\
                  S(\rho)=ns}}
 \frac1n S\!\left(\mathcal E_{\eta,N_B}^{\otimes n}(\rho)\right)
 =g\!\left(\eta g^{-1}(s)+(1-\eta)N_B\right).
\end{equation}
It is attained by the product thermal input \(\tau_{g^{-1}(s)}\). No independence assumption is imposed among the modes of \(\rho\); only the environment is required to be independent of that input.
\end{corollary}
\begin{proof}
The environment \(\tau_{N_B}\) has finite total mean photon number \(nN_B\) and entropy \(ng(N_B)\). The retained output has finite energy because passive mixing preserves combined total number, so its entropy is finite by Lemma~\ref{lem:energy-compactness}. Apply Theorem~\ref{thm:epni} to the independent inputs \(\rho\) and \(\tau_{N_B}\), and then apply the increasing function \(g\), to obtain \eqref{eq:thermal-attenuator-bound}.

For \(r=g^{-1}(s)<\infty\), the state \(\tau_r\) has finite total mean photon number \(nr\) and entropy \(ns\). Proposition~\ref{prop:thermal-mixing} gives
\[
 \mathcal E_{\eta,N_B}^{\otimes n}(\tau_r)
 =\tau_{\eta r+(1-\eta)N_B},
\]
whose entropy is \(n g(\eta r+(1-\eta)N_B)\). This attains the lower bound and proves \eqref{eq:thermal-attenuator-minimum}. The endpoints are included: at \(\eta=0\) the retained state is the environment, and at \(\eta=1\) it is the input. When \(s=0\) or \(N_B=0\), the corresponding thermal state is the vacuum, covered by the zero-mean case of Proposition~\ref{prop:thermal-mixing}.
\end{proof}

Corollary~\ref{cor:thermal-attenuator} concerns identical passive thermal attenuators with an independent product thermal environment. It asserts attainment, not uniqueness, and makes no claim for mode-dependent attenuations, amplifiers, or additive-noise channels.

The vacuum case also determines the classical capacity region of a
degraded pure-loss broadcast channel. Here one input mode per use is
passively split with auxiliary vacuum modes independent of the input
between receivers \(B,C\) and an inaccessible loss output. The power
transmissivities satisfy
\(0\le\eta_C<\eta_B\) and \(\eta_B+\eta_C\le1\), so \(C\)'s marginal
is obtained from \(B\)'s by a further pure-loss attenuation.
Section~\ref{sec:broadcast} gives the exact region for independent
classical messages and separate collective decoding, without feedback
or receiver cooperation. Its finite mean photon constraint is averaged
over the codebook and uses, and its finite-energy codewords may be
entangled across uses. The vacuum-port specialization of
Theorem~\ref{thm:epni} supplies the entropy input assumed in the
converse of Guha, Shapiro, and Erkmen~\cite{GSE2007};
coherent-state superposition coding supplies achievability.

\subsection{Context and earlier results}
Shannon introduced entropy power in his theory of communication
\cite{Shannon1948}; the Fisher-information proofs of Stam and Blachman
\cite{Stam1959,Blachman1965} became central to its later development.
For independent classical random vectors, the entropy power inequality
compares the entropy of a sum with the entropies of its summands. Its
bosonic counterpart replaces addition by passive mixing of quantum modes.
The analogy is especially natural for thermal states: their mean photon
numbers add with the beam-splitter weights, and their entropies are given
by the same function \(g\) appearing in Theorem~\ref{thm:epni}.
The formulation of Guha, Erkmen, and Shapiro
\cite[Section~II.B]{GES2007} makes this thermal comparison the central
conjecture and explains its consequences for bosonic communication.
Earlier work had established the capacity of the pure-loss bosonic channel
\cite{GGLMSY2004} and connected noisy-channel capacity questions with
minimum output entropy \cite{GGLMS2004}. The constrained minimum-output
problem also appears in the bosonic broadcast-channel analysis of Guha,
Shapiro, and Erkmen \cite{GSE2007}. Early EPnI special cases included
one-mode number-diagonal inputs with two nonzero probabilities and a
vacuum second port \cite{DSM2011}.

K\"onig and Smith \cite{KS2014} established a quantum entropy power
inequality using quantum Fisher information and a heat evolution. Their
results include the linear bound
\[
 S(\rho_C)\ge\eta S(\rho_A)+(1-\eta)S(\rho_B)
\]
for every transmissivity, and the exponential entropy power bound for a
balanced beam splitter. De Palma, Mari, and Giovannetti \cite{DMG2014}
proved the exponential bound for every transmissivity,
\begin{equation}\label{eq:exponential-qepi}
 e^{S(\rho_C)/n}
 \ge\eta e^{S(\rho_A)/n}+(1-\eta)e^{S(\rho_B)/n}.
\end{equation}
These results are formulated for an arbitrary finite number of modes and independent
inputs, with no product assumption within a group. De Palma, Mari, Lloyd, and Giovannetti
\cite{DMLG2015} subsequently treated more general linear mixing of
several independent groups of modes.
De Palma and Trevisan later gave a rigorous finite-energy treatment of the
Fisher-information and heat-flow steps using an integral formulation;
their conditional inequality with a trivial memory system yields
\eqref{eq:exponential-qepi} for arbitrary independent finite-energy inputs
\cite[Theorem~6]{DePalmaTrevisan2018}.

The entropy photon-number inequality requires a sharper comparison than
\eqref{eq:exponential-qepi} when the input entropy levels differ. When
$S(\rho_A)=S(\rho_B)$, the linear entropy bound already gives the EPnI
value by monotonicity of $g^{-1}$. For unequal entropy levels, thermal
inputs attain EPnI, whereas the exponential bound does not recover their
exact output entropy. This distinction matters for the
minimum output entropy problem with an input-entropy constraint: the
bound must reproduce the exact thermal comparison at unequal entropy
levels.

The EPnI itself was also known for two independent Gaussian inputs in any
finite number of modes, including Gaussian correlations within either input;
De Palma proves this case using symplectic eigenvalue comparisons
\cite[Appendix~A.8]{DePalmaThesis2017}.

The Gaussian optimizer theorem of Giovannetti, Holevo, and
Garc\'ia-Patr\'on \cite{GHG2015} settled the unconstrained minimum-output
problem for broad classes of Gaussian channels. Fixing a positive input
entropy is a different optimization problem. For one-mode gauge-covariant
channels, De Palma, Trevisan, and Giovannetti established optimality of
passive rearrangement at fixed spectrum \cite{DTG2016passive}, then the
sharp entropy bound for the quantum-limited attenuator
\cite{DTG2017attenuator}, and finally the sharp bound for one-mode
phase-covariant Gaussian channels with thermal noise
\cite[Theorem~4]{DTG2017}. Their multimode extension
\cite[Theorem~11]{DTG2017multimode} covers inputs diagonal in some
product basis, including classical correlations in that basis.

Later advances address further parts of this landscape. De Palma
\cite{DP2019} proved the sharp constrained bound for multimode
entanglement-breaking Gaussian attenuators and amplifiers with thermal
environments, and
improved lower bounds for other parameter regimes. Beigi and
Rahimi-Keshari \cite{BeigiRahimiKeshari2025} proved a one-mode
meta-logarithmic-Sobolev inequality that also gives a variational proof of
the one-mode constrained minimum-output result. Beigi and Mehrabi
\cite{BeigiMehrabi2026} established subset-convolution inequalities for
exponential entropy power and entropy monotonicity under repeated
symmetric quantum convolution of an identical input.

Theorem~\ref{thm:epni} supplies the thermal comparison for two freely
varying inputs, including arbitrary entanglement within either port.
Corollary~\ref{cor:thermal-attenuator} specializes it to the exact
finite-mode constrained minimum for identical thermal attenuators,
without a product or separability restriction on the finite-energy input.
Thermal states enter the proof of Theorem~\ref{thm:epni} as reference states and as limits of
regularized minimizers. The original inputs need not be Gaussian,
number-diagonal, or independent across modes.

\subsection{Proof strategy}
The difficulty is to compare the two freely varying inputs with thermal
states at their respective entropy levels. We fix those thermal reference
parameters near a hypothetical strictly violating pair. This turns the
contradiction into a negative minimum of a linear combination of three
entropies, after adding penalties that ensure compactness and regularity.
The reference parameters remain fixed throughout the variational argument.

An independent analytic ingredient is an interpolation theorem for positive diagonal quadratic forms, proved in Section~\ref{sec:interpolation}. For positive scalars $m$ and real parameters $x,y$, it transforms four comparisons with weights
\[
 mf(x)e^{x},\quad mf(x)e^{-x},\quad
 mf(y)e^{y},\quad mf(y)e^{-y},
 \qquad f(x)=\frac{x}{\sinh x},\quad f(0)=1,
\]
into the comparison with weight \(mf(x)f(y)/f(x+y)\). No compatibility is required between the linear maps and the diagonal multipliers. A Stieltjes representation for the derivative of an implicitly defined function is the essential analytic step. The theorem is formulated for finitely many spectral parameter values on source tests and permits unrestricted nonnegative target sums, so it also applies to the infinite matrix spaces used below.

Section~\ref{sec:minimization} develops the variational argument. Assuming a strict violation, we add a thermal port, a small thermal replacement of the output, a relative-entropy penalty, and a fourth-moment penalty. The resulting functional has faithful Gibbs minimizers with enough moment control for all subsequent traces. Varying one input at a time yields a component contraction in the logarithmic-mean metric \emph{at those minimizers}. It is not asserted as a universal channel inequality.

Section~\ref{sec:metric} combines two product-port variance estimates with the component contractions. These supply precisely the four hypotheses of the interpolation theorem. Its conclusion controls the norm of a centered thermal-balance defect: a linear functional measuring the creation--annihilation balance after subtracting the state's mean. In Section~\ref{sec:stationarity}, pairing the Gibbs identities with the relaxation generators gives an exact cancellation of the centered energy terms. As the fourth-moment penalty vanishes, the remaining defects and means tend to zero; their number-basis identities force both limiting inputs to be thermal. Entropy continuity then contradicts the original strictly negative objective value.

The auxiliary regularizations have linked roles: the relative-entropy penalty gives energy coercivity and supplies the strict margin in the final defect and mean estimate, while the thermal port supplies the slack needed for the associated Hessian comparisons. Output replacement controls the output logarithm, and the moment penalty justifies the unbounded-generator calculations. Only the last penalty tends to zero in the concluding compactness argument. This avoids assuming differentiability of entropy along an unbounded evolution or convergence of the input energies.

Section~\ref{sec:preliminaries} supplies the entropy, compactness, and thermal-mixing facts used in this argument. Each of the specialized interpolation, minimization, and limiting statements is proved within the paper.
